\documentclass[11pt, a4paper]{article}
\usepackage{amsmath,amssymb,amsfonts,amsthm}
\usepackage{geometry,hyperref,setspace,titlesec,booktabs,tabularx}
\geometry{left=2.5cm,right=2.5cm,top=2.5cm,bottom=2.5cm}
\setlength{\parindent}{0pt}\setlength{\parskip}{0.8em}\onehalfspacing
\hypersetup{colorlinks=true,linkcolor=blue,citecolor=blue,urlcolor=blue}
\newtheorem{proposition}{Proposition}
\title{\textbf{Memory as Propagation Terrain: \\ A Computational Theory Inspired by Neural Fields, Replay, Plasticity, and Free Energy}}
\author{Qing Zhou}\date{}
\begin{document}\maketitle

\begin{abstract}

This paper proposes a computational theory in which memory is understood as propagation terrain. Rather than treating memory primarily as static content that can be retrieved item by item, the model defines memory as the structural effect of past experience on the future propagation, competition, and convergence of activity. Once sensory input or an internal cue activates neural populations, activity does not travel through a blank network. It evolves through a state-dependent terrain jointly constituted by synaptic connectivity, excitation-inhibition balance, goal-dependent gating, value modulation, and replay traces. Repeatedly successful trajectories reduce the effective resistance of corresponding transitions; failed or dangerous trajectories strengthen risk-sensitive inhibition and gating; hippocampal-like replay converts a single experience into multiple offline updates; and free energy supplies a normative objective balancing prediction mismatch, complexity, conflict, energetic cost, and action value.

The paper develops a dual-timescale model. Fast dynamics describe activity propagation and attractor convergence in a gated excitatory-inhibitory neural field. Slow dynamics describe the formation of propagation terrain through locally eligible, value-modulated, and replay-amplified plasticity. The model further introduces a memory-field potential that represents the collective directional influence of many weak traces on candidate actions. This yields a hard criterion distinguishing the proposal from top-k retrieval: even when no individual memory is strong enough to change a decision, many directionally aligned but semantically dispersed traces should be able to produce a behavioral shift through aggregation.

The framework is intended both as a middle-level account of neural memory dynamics and as an implementable architecture for artificial agents with auditable memory, risk gating, and offline replay. It does not claim that the brain literally obeys a physical field law, nor that free energy, neural waves, hippocampal replay, or brain-region functions are interchangeable. It is a computable, simulable, comparable, and falsifiable integrative model.

\end{abstract}
\textbf{Keywords:} memory dynamics; propagation terrain; neural fields; hippocampal replay; synaptic plasticity; excitation-inhibition balance; free energy; active inference; agent memory; weak-evidence aggregation


\section{1. What Does Memory Preserve?}

Describing memory as storage is natural in engineering. An experience is encoded as an object, written to a location, and later retrieved by a cue. This description is useful for databases, but it does not fully explain several basic properties of biological memory:

\begin{enumerate}
\item Past experience changes present perception, attention, and action even when it is not explicitly recalled.
\item Repeated experiences form skills, habits, vigilance, and avoidance rather than merely adding records.
\item The same cue can activate different outcomes under different goals, bodily states, and contexts.
\item Memory structure continues to change during sleep and rest, when no new external event is occurring.
\item Behavioral influence often arises from many weak traces acting together rather than from one maximally similar episode.

\end{enumerate}
These observations suggest that memory preserves not only what happened, but also how future activity should flow. Experience changes connection strengths, local excitability, inhibitory boundaries, transition probabilities, value estimates, and control thresholds. When a related input later arrives, activity becomes more likely to spread in some directions, less likely to enter other states, and faster to converge on previously useful interpretations or actions.

The central thesis is therefore:

\begin{quote}
Memory is the shaping of future activity propagation and state transition by past experience.
\end{quote}

This shaping is called \textbf{propagation terrain}. It is not a mysterious additional medium in the brain. It is a middle-level description of the joint effects of connectivity, plasticity, excitation-inhibition balance, gating, and value modulation. A \textbf{memory field} is the distributed bias that this terrain exerts on current states and candidate actions.

Prediction, on this view, need not always mean that a system explicitly computes a picture of the future. Prediction can also appear when historically formed terrain makes present activity more likely to enter some low-cost states than others. The system behaves as if it anticipated the future partly because the past has already altered the routes available in the present.


\section{2. Theoretical Position: A Middle-Level Computational Model}

The framework sits between three levels of description.

The first is biological implementation. Real nervous systems consist of diverse cell types, synapses, dendrites, neuromodulators, oscillations, anatomical pathways, glia, bodily feedback, and developmental constraints. Wilson-Cowan population dynamics, continuous attractors and neural fields, timing-dependent plasticity, hippocampal sharp-wave ripples, and sequence replay provide partial mechanistic foundations.

The second is computational mechanism. At this level, the paper introduces propagation terrain, effective resistance, replay-mediated updating, gating, field potential, and weak-evidence aggregation. The question is which state variables are sufficient to generate the target behavior and how those variables interact.

The third is normative explanation. Variational and expected free energy provide a language for evaluating prediction mismatch, complexity, action risk, and preference satisfaction. The model does not assume that an individual neuron reads and minimizes a global free-energy scalar. Free energy is a normative description used by the modeler to evaluate the aggregate adaptive direction of local processes.

The framework therefore rejects the following literal equivalences:

\begin{itemize}
\item A memory field is not an independent physical field.
\item A neural activity wave is not memory content itself.
\item The hippocampus is not literally a computer cache or a capacitor.
\item Free energy is not identical to prediction error alone.
\item The prefrontal cortex is not a unique central executive.
\item The brain is not assumed to solve the equations below exactly.

\end{itemize}
The value of these concepts lies not in literal identity, but in whether they compose into a model that can be simulated, compared with alternatives, and rejected by evidence.


\section{3. Relation to Existing Frameworks}

\subsection{3.1 Neural Fields and Population-Level Propagation}

Neural-field theory describes large-scale neural dynamics through population activity distributed over a continuous representational or anatomical space. Local activity is influenced by nearby and distant populations, conduction delays, nonlinear response functions, and external input. Depending on the connectivity kernel and time constants, neural fields can produce stationary patterns, traveling waves, rotating waves, oscillations, and attractors.

The present theory inherits this population-dynamical language but adds a central assumption: the connectivity kernel is not merely a fixed background. It is continually shaped by experience, value, replay, and goal-dependent gating. The field therefore propagates activity while gradually altering the future conditions of propagation.

\subsection{3.2 Excitation-Inhibition Balance}

Wilson-Cowan models demonstrate that coupled excitatory and inhibitory populations can generate rich collective dynamics. The present model explicitly distinguishes excitatory and inhibitory populations rather than allowing a single synapse to switch arbitrarily between positive and negative identity.

In this interpretation, a practiced trajectory can appear as increased effective excitatory gain for selected transitions. Risk learning can appear as strengthened inhibitory boundaries, higher gating thresholds, or improved access to avoidance states. Together, excitation and inhibition determine whether activity can enter, cross, or leave a region of state space.

\subsection{3.3 Plasticity and Complementary Learning Systems}

Synaptic plasticity allows co-activity, temporal order, error signals, and neuromodulation to change future connectivity. Complementary learning systems theory further distinguishes a rapidly learning hippocampal system from a slowly integrating cortical system. The fast system preserves specific episodes; the slow system gradually extracts statistical structure while limiting catastrophic interference.

The present model expresses this division through two learning timescales. Fast traces bind new trajectories. Offline replay selectively resamples them. Slow connectivity then accumulates repeated updates into stable priors and reusable transition structure.

\subsection{3.4 Replay, Reactivation, and Offline Learning}

Hippocampal sequences can be reactivated during rest, sleep, and pauses in decision making. Replay is not necessarily a literal recording of past experience. It can be temporally compressed, reversed, reordered, and biased by reward, novelty, and current goals.

The present theory treats replay as an \textbf{endogenous training-distribution generator}. Replay determines which past trajectories receive additional learning iterations. Whether one experience becomes a stable route depends not only on its online intensity, but also on whether it is replayed, in what order, and under which value and risk signals.

\subsection{3.5 Free Energy, Predictive Processing, and Active Inference}

The free-energy principle and active inference describe perception, learning, and action as processes that reduce uncertainty while maintaining preferred states. The present framework uses this language conservatively:

\begin{quote}
Free energy characterizes system-level mismatch and adaptive cost; it is not assumed to be a biological quantity explicitly read by each local circuit.
\end{quote}

The main contribution is not another statement that a system minimizes free energy. It is a concrete middle-level candidate mechanism: activity evolves through propagation terrain; replay and local plasticity alter that terrain; gating reconfigures it according to goals and conflict; and free energy evaluates whether these changes reduce expected future adaptive cost.


\section{4. Core Concepts: Field, Wave, Path, and Resistance}

\subsection{4.1 Memory Trace}

Let an episodic trace be

\[
m_i=(c_i,a_i,o_i,v_i,r_i,t_i,\gamma_i),
\]

where $c_i$ is context, $a_i$ is an action or transition, $o_i$ is outcome, $v_i$ is value, $r_i$ is risk or failure, $t_i$ is time, and $\gamma_i$ is confidence or precision.

A trace may be explicitly represented or distributed across connection changes. It is a local record of an experience, not the whole memory field.

\subsection{4.2 Propagation Terrain}

Define the middle-scale propagation terrain as

\[
\mathcal{T}_t=\{W_t^{XY},G_t^{XY},\Theta_t^X,H_t,\mathcal{B}_t\},\qquad X,Y\in\{E,I\},
\]

where:

\begin{itemize}
\item $W_t^{XY}$ denotes cell-type-specific connectivity kernels;
\item $G_t^{XY}$ denotes goal-, attention-, and context-sensitive gating;
\item $\Theta_t^X$ denotes adaptive thresholds or local excitability;
\item $H_t$ denotes fast episodic traces and replay buffers;
\item $\mathcal{B}_t$ denotes modulatory variables including value, risk, and bodily state.

\end{itemize}
Propagation terrain is a state collection rather than a single matrix. It determines the ease with which current activity can move between regions of state space.

\subsection{4.3 Memory Wave}

A memory wave is the spatiotemporal evolution of activity through propagation terrain:

\[
u_t(x)\longrightarrow u_{t+\Delta t}(x).
\]

It may appear as local diffusion, sequence propagation, a traveling wave, attractor convergence, or competitive extinction. A memory wave is not an entity independent of neural organization. It is a description of system-state evolution.

\subsection{4.4 Effective Propagation Resistance}

For an operational engineering measure, define the effective conductance from state region $i$ to $j$ as

\[
K_{ij}^{\mathrm{eff}}(t)=G_{ij}(t)W_{ij}(t)S_{ij}(t),
\]

where $S_{ij}$ is the match between the current state and the route. The corresponding effective resistance is

\[
R_{ij}^{\mathrm{eff}}(t)=\frac{1}{\epsilon+\left|K_{ij}^{\mathrm{eff}}(t)\right|}.
\]

This resistance is not a literal electrical resistance. It operationalizes the stimulus intensity, time, energy, or control investment required for a transition. If learning allows a related input to reach a target state with less perturbation or in less time, the effective resistance of that route has decreased.

\subsection{4.5 Memory-Field Potential}

For a current context $q$ and candidate action $a$, define

\[
\Phi_M(a|q)=\sum_i\gamma_i d_i K(q,a,m_i),
\]

where $K(q,a,m_i)$ measures contextual transfer, $\gamma_i$ measures trace strength and precision, and $d_i\in[-1,1]$ indicates whether the trace supports or inhibits the action.

The essential feature is not the largest term, but the ability of many small terms to accumulate. Action selection can be written as

\[
a^\ast=\arg\min_a\left[\mathcal{L}_{task}(a|q)+\lambda_r\operatorname{Risk}(a|q)+\lambda_c\operatorname{Conflict}(a|q)-\lambda_m\Phi_M(a|q)\right].
\]

A semantically dispersed set of failures can therefore jointly favor checking before acting. Likewise, individually weak successful episodes can collectively form a stable skill bias.


\section{5. Fast Timescale: A Gated Excitatory-Inhibitory Neural Field}

Let $u_E(x,t)$ and $u_I(x,t)$ denote excitatory and inhibitory population activity over an anatomical or representational space $x$. A simplified gated neural field is

\[
\begin{aligned}
\tau_E\frac{\partial u_E(x,t)}{\partial t}
&=-u_E(x,t)+\phi_E\left[
\int_\Omega G_{EE}(x,y,t)W_{EE}(x,y,t)u_E(y,t)\,dy\right.\\
&\qquad\left.-\int_\Omega G_{EI}(x,y,t)W_{EI}(x,y,t)u_I(y,t)\,dy
+I_E(x,t)-\Theta_E(x,t)\right],\\
\tau_I\frac{\partial u_I(x,t)}{\partial t}
&=-u_I(x,t)+\phi_I\left[
\int_\Omega G_{IE}(x,y,t)W_{IE}(x,y,t)u_E(y,t)\,dy\right.\\
&\qquad\left.-\int_\Omega G_{II}(x,y,t)W_{II}(x,y,t)u_I(y,t)\,dy
+I_I(x,t)-\Theta_I(x,t)\right].
\end{aligned}
\]

Here $\phi_E$ and $\phi_I$ are population response functions, $W_{XY}$ are experience-shaped connectivity kernels, $G_{XY}$ are task- and context-dependent gates, $\Theta_X$ are adaptive thresholds, and $I_X$ contains sensory or internally replayed input.

The equation separates three sources of constraint. Connectivity $W$ preserves slow statistical structure and determines which states tend to co-activate. Gating $G$ temporarily opens or closes routes according to the present goal, allowing the same long-term terrain to be reconfigured for different tasks. Thresholds and inhibition determine when a region can be entered and prevent all strong routes from spreading indiscriminately.

Propagation terrain is therefore not a fixed map, but a conditional map jointly generated by structural priors and current control:

\[
\mathcal{T}_t=\mathcal{T}(W_t,G_t,\Theta_t,u_I(t),\mathcal{B}_t).
\]

This explains why the same memory cue can produce different outcomes in different states. Long-term connectivity supplies possible routes; current gating determines which possibilities are active now.


\section{6. Slow Timescale: Value- and Replay-Modulated Plasticity}

Activity propagation becomes memory formation only when it changes future propagation conditions. Let $e_{ij}^{XY}(t)$ be a local eligibility trace generated by co-activity and temporal relation:

\[
\tau_e\dot e_{ij}^{XY}(t)=-e_{ij}^{XY}(t)+\psi_X(u_i^X(t))\psi_Y(u_j^Y(t-\delta)).
\]

Eligibility only indicates that a connection recently participated. It does not determine the direction of learning by itself. Connectivity is updated by local eligibility, modulatory signals, homeostatic constraints, and replay:

\[
\dot W_{ij}^{XY}
=\eta_{XY}\delta_t^{XY}e_{ij}^{XY}
-\beta_{XY}\left(W_{ij}^{XY}-W_{0,ij}^{XY}\right)
-\lambda_{homeo}\nabla_{W_{ij}^{XY}}\mathcal{H}
+\rho\mathcal{R}_{ij}^{XY}(t).
\]

Here $\delta_t^{XY}$ is a value, prediction-error, novelty, or risk signal; $\beta_{XY}$ describes forgetting or return to baseline; $\mathcal{H}$ is a homeostatic constraint; $\mathcal{R}_{ij}^{XY}$ is replay-induced updating; and $\rho$ is replay gain.

Successful experience should not be interpreted as strengthening every participating connection. A more defensible claim is that transitions that reduce future task cost while satisfying homeostatic constraints tend to gain effective conductance. Failure does not simply weaken a trajectory. It may strengthen access to risk representations, reinforce inhibitory boundaries, or raise the gate threshold for a dangerous action.

Thus,

\[
\text{successful experience}\Rightarrow\text{effective resistance of useful routes tends to decrease},
\]

and

\[
\text{failed experience}\Rightarrow\text{risk detection and inhibitory gating tend to strengthen}.
\]

Learning not to touch fire does not erase fire-related memory. It makes the cue propagate more reliably toward risk recognition and avoidance.


\section{7. Replay: From One Experience to Multiple Internal Updates}

\subsection{7.1 Replay as Endogenous Input}

During online interaction, external input drives the field:

\[
I_X(x,t)=I_X^{\mathrm{sensory}}(x,t).
\]

During offline or intermittent periods, a fast memory system generates replay input:

\[
I_X(x,t)=I_X^{\mathrm{replay}}(x,t;z_k),
\]

where $z_k$ is a selected event trajectory. A replayed trajectory may preserve, reverse, compress, or reorganize the original sequence. It reactivates eligibility structure and creates additional opportunities to update slow connectivity.

\subsection{7.2 Replay Selection}

Not every event should be replayed with equal probability. Define replay priority as

\[
P_{\mathrm{replay}}(z_k)\propto
\exp\left[
\alpha_NN_k+\alpha_V|V_k|+\alpha_UU_k+\alpha_RR_k-\alpha_CC_k
\right],
\]

where $N_k$ is novelty, $|V_k|$ is value magnitude, $U_k$ is uncertainty or unresolved prediction error, $R_k$ is current-goal relevance, and $C_k$ is replay and update cost.

Replay is therefore not mechanical copying. It is selective retraining biased by unresolved problems, value, risk, and current goals.

\subsection{7.3 The Dual Role of Replay}

Replay can have at least two potentially opposing effects.

First, consolidation: repeated reactivation of valuable trajectories can stabilize slow cortical-like connectivity and reduce convergence time for related future inputs.

Second, revision: when a trajectory conflicts with new evidence, replay can place the old route into a new context, weakening it, splitting it, or attaching conditional gates.

Strong replay therefore does not necessarily mean stronger verbatim preservation. Replay is another training opportunity, and its result depends on the error, value, and context present during replay.


\section{8. Free Energy as a Normative Objective, Not a Neural Magic Variable}

Let the system maintain an approximate posterior $q(s)$ over hidden states and a generative model $p(o,s|M)$. Variational free energy is

\[
F[q,M]=\mathbb{E}_{q(s)}[\ln q(s)-\ln p(o,s|M)].
\]

It can be expressed as a complexity-accuracy tradeoff:

\[
F=D_{KL}[q(s)\Vert p(s|M)]-\mathbb{E}_{q(s)}[\ln p(o|s,M)].
\]

For operational purposes, define a system cost

\[
\mathcal{J}_t
=F_t+\lambda_EE_t+\lambda_{\Delta W}\|\dot W_t\|^2
+\lambda_{\mathrm{conf}}\operatorname{Conf}_t+\lambda_G\|G_t\|^2-\lambda_VV_t.
\]

The terms represent present model mismatch and complexity, activity cost, rapid rewriting cost, path conflict, control cost, and preference or action value.

Gating can be modeled as approximate descent on expected cumulative cost:

\[
\dot G_t=-\kappa_G\nabla_G\mathbb{E}\left[\int_t^{t+H}\mathcal{J}_\tau d\tau\right]+\xi_G(t).
\]

Slow connectivity is modified by local plasticity and replay in a direction that approximately improves future cost:

\[
\dot W_t\approx-\kappa_W\widehat{\nabla_W\mathbb{E}[\mathcal{J}_{future}]}.
\]

The hat emphasizes that a biological system need not obtain an exact global gradient. Local eligibility and modulatory signals approximate it only under particular conditions.

The relation between propagation terrain and free energy can be summarized as

\[
\begin{aligned}
\text{terrain matches present input}
&\Rightarrow \text{faster convergence and less conflict}\\
&\Rightarrow \text{less control demand}\\
&\Rightarrow \mathcal{J}\text{ tends to decrease}.
\end{aligned}
\]

A low-resistance route is not necessarily a good route. Habit, addiction, bias, and compulsive behavior may all be locally easy but globally costly. The role of value and expected cost is to distinguish mere repetition of an easy route from movement toward a viable long-term state.


\section{9. Unified Dual-Timescale Model}

Activity propagation, local plasticity, replay, and gating can be summarized as

\[
\begin{aligned}
\tau_E\dot u_E
&=-u_E+\phi_E\left(G_{EE}\odot W_{EE}u_E-G_{EI}\odot W_{EI}u_I+I_E-\Theta_E\right),\\
\tau_I\dot u_I
&=-u_I+\phi_I\left(G_{IE}\odot W_{IE}u_E-G_{II}\odot W_{II}u_I+I_I-\Theta_I\right),\\
\tau_e\dot e_{ij}^{XY}
&=-e_{ij}^{XY}+\psi_X(u_i^X)\psi_Y(u_j^Y),\\
\dot W_{ij}^{XY}
&=\eta_{XY}\delta_t^{XY}e_{ij}^{XY}
-\beta_{XY}(W_{ij}^{XY}-W_{0,ij}^{XY})
-\lambda_h\nabla_W\mathcal{H}
+\rho\mathcal{R}_{ij}^{XY},\\
\dot G
&=-\kappa_G\nabla_G\mathbb{E}\left[\int_t^{t+H}\mathcal{J}_\tau d\tau\right]+\xi_G,\\
\mathcal{J}_t
&=F_t+\lambda_EE_t+\lambda_{\Delta W}\|\dot W_t\|^2
+\lambda_{\mathrm{conf}}\operatorname{Conf}_t+\lambda_G\|G_t\|^2-\lambda_VV_t.
\end{aligned}
\]

The model contains three broad timescales:

\begin{enumerate}
\item Milliseconds to seconds: activity propagation, competition, and gating.
\item Seconds to hours: eligibility, rapid binding, and replay selection.
\item Days and beyond: slow consolidation, forgetting, and structural reorganization.

\end{enumerate}
Its basic loop is

\[
\begin{aligned}
\text{input}
&\rightarrow\text{activity propagation}
\rightarrow\text{competition and action}\\
&\rightarrow\text{error, value, and risk feedback}
\rightarrow\text{fast trace}\\
&\rightarrow\text{selective replay}
\rightarrow\text{slow terrain update}\\
&\rightarrow\text{changed future propagation}.
\end{aligned}
\]

Memory is therefore not one isolated module in the loop. It is the loop's capacity to alter its own future dynamics.


\section{10. Four Propositions Derived from Propagation Terrain}

\subsection{Proposition 1: Repeated Success Shortens Convergence Time, with Saturating Returns}

If a class of inputs repeatedly follows the same route to a high-value state, and homeostatic constraints permit the effective conductance of that route to increase, then

\[
K_{ij}^{\mathrm{eff}}\uparrow
\Rightarrow R_{ij}^{\mathrm{eff}}\downarrow
\Rightarrow T_{\mathrm{conv}}\downarrow.
\]

Because homeostasis, adaptive thresholds, and bounded weights constrain growth, improvement should eventually saturate rather than accelerate without limit. Skill acquisition should therefore be rapid early and show diminishing returns later.

\subsection{Proposition 2: Failure Learning Can Accelerate Risk Recognition}

If failed trajectories are replayed and strengthen propagation from cues to risk representations, the system can recognize danger faster while becoming less likely to enter the failed action:

\[
T_{\mathrm{risk-detect}}\downarrow,\qquad P(a_{\mathrm{failed}}|q)\downarrow.
\]

Failure memory can therefore create both a high-conductance risk-recognition route and a high-resistance dangerous-action route.

\subsection{Proposition 3: Conflict Temporarily Increases Gating Investment}

When several routes are simultaneously active and similarly valued, conflict rises. If additional gating reduces future cumulative cost, then

\[
\operatorname{Conf}_t\uparrow\Rightarrow\|G_t-G_{\mathrm{base}}\|\uparrow.
\]

The system invests control to separate competing routes. If control becomes too expensive or resources are exhausted, however, gating may collapse and behavior may revert to habitual low-resistance routes. This yields a non-monotonic prediction: moderate conflict increases control, whereas extreme load can reduce it.

\subsection{Proposition 4: Weak Evidence Can Cross a Decision Threshold at the Field Level}

Suppose every individual contribution to action $a$ is subthreshold:

\[
|\gamma_i d_iK(q,a,m_i)|<\theta,\quad\forall i.
\]

If their directionally aligned sum satisfies

\[
\left|\sum_i\gamma_i d_iK(q,a,m_i)\right|\ge\theta,
\]

then field potential changes the action even though item-wise threshold retrieval would not. This is the most important discriminative proposition of the theory.


\section{11. Memory Fields Versus Top-k Retrieval}

A conventional retrieval-augmented system scores the similarity between a query and stored items, then returns the highest-scoring $k$ items:

\[
\mathcal{M}_{top-k}(q)=\operatorname{arg\,topk}_i\operatorname{sim}(q,m_i).
\]

This works well for directly relevant facts, but it can miss experiences that are semantically dissimilar while sharing the same action implication. Consider four failures:

\begin{itemize}
\item A deployment failed because environment variables were not checked.
\item A paper contained an error because citations were not verified.
\item An analysis had to be repeated because units were not checked.
\item An automation damaged files because the target path was not confirmed.

\end{itemize}
The objects differ and may not be near one another in embedding space. Yet all support the same action direction: \textbf{verify before producing a high-impact output}.

A memory-field system stores not only an event embedding, but also its directional effect on action primitives:

\[
m_i=(e_i,a_i^{support},a_i^{inhibit},v_i,r_i,\gamma_i).
\]

Its candidate-action potential can be written as

\[
\Phi_M(a|q)=\sum_i
\underbrace{K_c(q,e_i)}_{\text{context transfer}}
\cdot\underbrace{K_a(a,a_i)}_{\text{action direction}}
\cdot\underbrace{\gamma_i(v_i-r_i)}_{\text{value and risk}}.
\]

An implementation need not scan every memory. Action-primitive indexes, low-rank field representations, graph diffusion, kernel approximation, or hierarchical aggregation can approximate the potential. The approximation must, however, preserve the ability of many weak traces to jointly cross a threshold.

The hard test is:

\begin{quote}
Under a controlled computational budget, can a memory-field model reliably solve weak-evidence aggregation tasks that top-k similarity retrieval cannot?
\end{quote}

If not, propagation terrain adds terminology but no distinct computational capacity.


\section{12. An Implementable Architecture for Artificial Agents}

A minimal agent implementation contains six components.

\subsection{12.1 Event Encoder}

Each interaction is written as a structured event:

\[
m_i=(context,goal,action,outcome,value,risk,confidence,time).
\]

Text preserves detail; structured fields enable aggregation across tasks.

\subsection{12.2 Fast Episodic Buffer}

New events first enter a fast buffer rather than immediately rewriting long-term policy. The buffer performs rapid binding and enables post-task replay.

\subsection{12.3 Action-Primitive Graph}

The system maintains reusable action primitives such as verify first, request clarification, search for evidence, make the minimum change, roll back, or stop a high-risk operation. Events create support and inhibition edges on these primitives through success, failure, and risk signals.

\subsection{12.4 Memory-Field Aggregator}

For the current context and candidate action, compute

\[
\Phi_M(a|q)=\Phi_{episodic}(a|q)+\Phi_{semantic}(a|q)+\Phi_{risk}(a|q)+\Phi_{goal}(a|q).
\]

The output is not merely retrieved text. It is a directional bias over actions together with an auditable decomposition of its evidence.

\subsection{12.5 Offline Replay Engine}

At task completion, idle time, or after an error, replay selects trajectories according to novelty, value, risk, and unresolved mismatch. It can generate compressed summaries of successful trajectories, counterfactual replays of failure, clusters of cross-task regularities, conflict checks for obsolete rules, and proposed updates to action-primitive weights.

\subsection{12.6 Audit and Homeostatic Constraints}

To prevent one intense event or erroneous replay from permanently distorting the agent, long-term updates should be bounded, source-traceable, defeasible by conflicting evidence, context-conditional, reversible under counterevidence, and subject to higher confidence requirements when their impact is high.

The fast buffer and replay engine, slow structural aggregator, and goal-risk gating are functionally analogous to selected aspects of hippocampal, cortical, and prefrontal organization. They are engineering mappings, not one-to-one copies of brain regions.


\section{13. Simulation and Experimental Designs}

\subsection{13.1 Neural-Dynamical Simulation}

A two-dimensional or graph-structured neural field can contain two candidate routes leading to different attractors. Training repeatedly presents trajectories with distinct value feedback. Compare fixed connectivity, online plasticity without replay, online plasticity with uniform replay, online plasticity with value-selective replay, and the full model with goal gating and conflict cost.

Primary measures include attractor convergence time, input intensity required to reach a target, route error rate, transfer to new tasks, retention after interference, latency from risk cue to avoidance, total activity, and gating cost.

The model predicts that value-selective replay will consolidate valuable routes more efficiently than uniform replay, but that excessive selection can also produce bias and route rigidity.

\subsection{13.2 Weak-Evidence Aggregation Benchmark}

Construct tasks in which every past event has low semantic similarity to the present task, yet the events jointly support one action principle. Compare top-k vector retrieval, full-history summarization, rule memory, memory-field aggregation, a model without action-direction labels, and a model without replay.

Success is not remembering more text. It is choosing the action jointly supported by weak evidence without a large increase in false positives.

\subsection{13.3 Counterfactual Replay Experiment}

After one failed trajectory, compare no replay, verbatim replay, reverse replay, and counterfactual replay containing an alternative action. Test later risk recognition and transfer to new contexts. The model predicts that verbatim replay mainly strengthens episode memory, whereas counterfactual replay more effectively creates transferable action gating.

\subsection{13.4 Biological and Behavioral Predictions}

If propagation terrain is explanatory, then:

\begin{enumerate}
\item Replay strength and related cortical representational change should jointly predict later convergence speed, not only explicit recall.
\item Failure learning should be able to improve risk-cue detection while suppressing dangerous action.
\item Repetition should reduce activation threshold and response time, with saturating improvement.
\item Moderate conflict should increase control-related signals, whereas extreme load may trigger habitual responding.
\item Multiple individually subthreshold but directionally aligned cues should jointly alter choice.
\item Disrupting offline replay should impair cross-context route consolidation, not only episodic detail.

\end{enumerate}
These predictions are not unique to this model. A serious test requires quantitative comparison with alternatives.


\section{14. Falsification Conditions}

A theoretical model has content only if it can fail. The propagation-terrain framework faces at least six falsification conditions.

First, if many weak memories do not produce reliable behavioral effects beyond top-k retrieval after similarity, value, frequency, and computational budget are controlled, memory-field potential has no distinct computational value.

Second, if replay improves only explicit episodic recall but does not change convergence time, transfer, or gating structure, the claim that replay shapes propagation terrain must be narrowed.

Third, if failure produces only generalized suppression and never the joint effect of faster risk recognition and reduced dangerous action, the proposed risk-gating mechanism is inadequate.

Fourth, if propagation-terrain variables do not explain behavior or neural data better than simpler fixed-attractor, reinforcement-learning, or retrieval models, the simpler models should be preferred.

Fifth, if arbitrary post-hoc adjustment of $W$, $G$, value, and free-energy weights can explain any result, the framework is unfalsifiable. Update rules, parameter ranges, and evaluation criteria must be preregistered in empirical tests.

Sixth, if replay selection only reinforces existing high-value routes and cannot be revised by new evidence, the model produces rigidity and confirmation bias. Without homeostasis, exploration, and counterevidence, it should not be considered an adaptive memory architecture.


\section{15. Limitations and Boundaries}

This is an integrative middle-level theory, not a complete neuroscience model. It omits neuronal diversity, dendritic computation, neuromodulatory detail, glia, anatomical constraints, interoception, and development.

Effective resistance, propagation terrain, and memory-field potential are operational concepts. They may be estimated through convergence time, activation threshold, transition probability, control cost, or action bias, but no unique measurement currently exists.

Neural fields are useful for population dynamics but may hide the importance of discrete symbols, compositional structure, and long-range communication. Real cognition probably combines continuous dynamics with discrete organization.

The free-energy framework is highly general. Unless a model specifies its generative assumptions, observations, preferences, update rules, and alternatives, free energy can become an explanation compatible with any result. The present theory therefore places falsifiability primarily on replay, convergence time, risk gating, and weak-evidence aggregation rather than treating free-energy reduction itself as sufficient evidence.

Low resistance is not equivalent to correctness. Habit, bias, addiction, and trauma responses can all create strong routes. Adaptation depends on value learning, environmental feedback, counterfactual replay, exploration, and meta-control, not merely on making familiar paths easier.

Finally, structured events, action primitives, and audit mechanisms in the agent implementation are more explicit than their biological analogues. They are neuroscience-inspired engineering designs, not literal claims about the brain.


\section{16. Definitions, Established Mechanisms, and Hypotheses}

\begin{center}\begin{tabularx}{\textwidth}{X p{0.24\textwidth}}\toprule
\textbf{Claim} & \textbf{Status} \\
\midrule
Propagation terrain is the collection of connectivity, gating, thresholds, replay traces, and modulatory states & Model definition \\
Effective resistance is operationalized as inverse effective conductance & Engineering definition \\
Memory-field potential is a weighted sum of directional traces & Model definition \\
A sum of subthreshold contributions can cross a decision threshold & Mathematical result \\
Neural fields can generate waves, oscillations, and attractors & Established model result \\
Local eligibility and modulatory signals can alter connectivity & Empirically supported mechanism class \\
Hippocampal-like replay accelerates consolidation of slow propagation terrain & Testable theoretical hypothesis \\
Successful experience tends to lower resistance along useful routes & Testable theoretical hypothesis \\
Failure tends to form risk-recognition routes and inhibitory gates & Testable theoretical hypothesis \\
Free energy is a sufficient unifying objective for these local processes & Normative hypothesis \\
The biological brain exactly solves the unified equations & Not claimed \\
Memory is an independent physical wave or field & Explicitly rejected \\
\bottomrule\end{tabularx}\end{center}


\section{17. Final Thesis}

The complete loop can be condensed as follows:

\begin{quote}
External input and internal cues evoke population activity. Activity evolves through propagation terrain jointly formed by historical connectivity, excitation-inhibition balance, and current gating. Action outcomes generate value, risk, and prediction-error signals. Local plasticity records recently involved routes. Replay converts selected experiences into additional internal training. Slow structure changes, making future related inputs more likely to enter some states and less likely to enter others. Free energy supplies a normative evaluation of whether this loop reduces long-term mismatch, conflict, energetic cost, and action risk.
\end{quote}

In one sentence:

\begin{quote}
Memory is not a static warehouse of past content, but the continuing shaping of future activity propagation, state-transition difficulty, and action selection by past experience.
\end{quote}

If this view is correct, the key capacity of memory is not merely to find the most similar past. It is to allow many dispersed, weak, and directionally aligned experiences to jointly alter the present. For neuroscience, this requires placing memory traces, activity propagation, offline replay, and goal-dependent gating in one dynamical loop. For artificial intelligence, it requires moving beyond top-k retrieval toward memory architectures that accumulate directional evidence, reorganize themselves offline, and remain auditable.


\section*{References}

Amari, S. (1977). Dynamics of pattern formation in lateral-inhibition type neural fields. \textit{Biological Cybernetics}, 27, 77-87.

Bi, G. Q., \& Poo, M. M. (1998). Synaptic modifications in cultured hippocampal neurons: Dependence on spike timing, synaptic strength, and postsynaptic cell type. \textit{Journal of Neuroscience}, 18(24), 10464-10472.

Buzsáki, G. (2015). Hippocampal sharp wave-ripple: A cognitive biomarker for episodic memory and planning. \textit{Hippocampus}, 25(10), 1073-1188.

Foster, D. J., \& Wilson, M. A. (2006). Reverse replay of behavioural sequences in hippocampal place cells during the awake state. \textit{Nature}, 440, 680-683.

Friston, K. (2010). The free-energy principle: A unified brain theory? \textit{Nature Reviews Neuroscience}, 11, 127-138.

McClelland, J. L., McNaughton, B. L., \& O'Reilly, R. C. (1995). Why there are complementary learning systems in the hippocampus and neocortex: Insights from the successes and failures of connectionist models of learning and memory. \textit{Psychological Review}, 102(3), 419-457.

O'Keefe, J., \& Recce, M. L. (1993). Phase relationship between hippocampal place units and the EEG theta rhythm. \textit{Hippocampus}, 3(3), 317-330.

Sutton, R. S. (1991). Dyna, an integrated architecture for learning, planning, and reacting. \textit{ACM SIGART Bulletin}, 2(4), 160-163.

Wilson, H. R., \& Cowan, J. D. (1972). Excitatory and inhibitory interactions in localized populations of model neurons. \textit{Biophysical Journal}, 12(1), 1-24.

Wilson, M. A., \& McNaughton, B. L. (1994). Reactivation of hippocampal ensemble memories during sleep. \textit{Science}, 265(5172), 676-679.
\end{document}
